\chapter{Die Oberfläche}
\label{ch:oberflaeche}

\section{Starten}

\begin{itemize}
  \item \textbf{In der Gesamt-GUI der Gruppe:} \pfad{start_gui.bat} im Repository-Ordner, Reiter
    \enquote{M4 Behavior Cloning}. Direkt in diesen Reiter: \pfad{start_gui.bat --tab m4}.
  \item \textbf{Eigenes Fenster:} \pfad{code\m4_behavior_cloning\BC-Oberflaeche.bat}.
\end{itemize}
Beides ist dieselbe Oberfläche mit denselben Reitern. In der Gesamt-GUI teilen sich Modul~2 und
Modul~4 den Roboter, Modul~1 und Modul~4 die Draufsichtkamera automatisch: Wer verbindet,
bekommt das Gerät, die andere Seite wird getrennt. Läuft auf der anderen Seite gerade etwas, kommt
stattdessen eine Meldung.

\section{Backend-Leiste}
\label{sec:backend}

Ganz oben, für alle Reiter gleich:
\begin{description}
  \item[Backend] \enquote{Simulation} (eigener SimRobot im Prozess) oder \enquote{Neura}
    (\acs{vm} oder Anlage).
  \item[Override] Geschwindigkeitsfaktor der Steuerung. Für Aufnahme und Fahrt \num{0.5}; das
    Modell erzwingt den Wert seiner Aufnahme.
  \item[Reale Anlage freigeben] Ohne diesen Haken verweigert der Adapter jede Bewegung, sobald
    die Steuerung nicht als Simulation antwortet.
  \item[\knopf{Verbinden} / \knopf{Trennen}] Verbindung zur Steuerung.
\end{description}

Die farbige Anzeige zeigt nicht die Auswahl, sondern die Antwort der Steuerung:
\begin{center}
  \small
  \begin{tabular}{@{}l l@{}}
    \toprule
    Farbe & Bedeutung \\
    \midrule
    grau & nicht verbunden \\
    blau & SimRobot im Prozess, keine Steuerung \\
    grün & Steuerung meldet Simulation (\acs{vm}) \\
    orange & Simulation nicht bestätigt, Bewegung gesperrt \\
    \textbf{rot} & \textbf{reale Anlage, Bewegung freigegeben} \\
    \bottomrule
  \end{tabular}
\end{center}

\section{Reiter Systemcheck}

Liste aller Prüfungen in drei Gruppen: \emph{hardwarefrei}, \emph{verbunden ohne Bewegung} und
\emph{mit Bewegung}. Einen Punkt auswählen zeigt, was er tut; \knopf{Diesen Punkt ausführen}
oder anhaken und \knopf{Angehakte ausführen}. Jeder Punkt endet mit
\texttt{[OK]}, \texttt{[WARNUNG]} oder \texttt{[FEHLER]}. \knopf{Protokoll speichern} legt das
Ergebnis unter \pfad{Berichte/} ab. \knopf{Auf Änderungen prüfen} und
\knopf{Oberfläche neu starten} laden geänderten Code. \Cref{tab:checks} in
\cref{ch:werkzeuge} beschreibt die Prüfungen.

\section{Reiter Aufnahme}

Von oben nach unten:
\begin{description}
  \item[1. Kameras] Plätze \texttt{wrist} und \texttt{scene} besetzen (Daheng per
    Seriennummer, Webcam oder Simulation). \knopf{Vorschau starten} zeigt beide Bilder;
    \knopf{Kamera \dots} öffnet die Kameraeinstellungen (\cref{sec:kamera-einst}).
  \item[2. Ablauf] Ablaufdatei wählen. Bearbeitet wird immer eine Arbeitskopie
    (\knopf{Aus Original zurücksetzen}). Die Punkte-Tabelle zeigt jeden Punkt mit seiner Lage;
    daneben \knopf{Touch-up} (aktuelle Armstellung übernehmen), \knopf{Touch-up rückgängig},
    \knopf{Teach-Serie (POINT am Arm)} und \knopf{$\blacktriangle$ Z hoch (halten, F7)}.
  \item[3. Aufnahme] Episoden, Rauschen, Tempo, Lagen und Varianten
    (\cref{sec:aufnahme-einst}), dann \knopf{Aufnahme starten}. Roter \knopf{STOPP} jederzeit.
  \item[4. Metadaten] Block, Objekt, Licht, Kameralage. Sie landen in jeder Episode.
\end{description}
Rechts eine Abdeckungskarte der bisher aufgenommenen Lagen (optional über alle Tage).

\section{Reiter Training}

\begin{description}
  \item[1. Datensatz] Aufnahmen auswählen (\knopf{Hinzufügen\dots}), Ziel angeben,
    \knopf{Nur prüfen} oder \knopf{Exportieren}. \enquote{nur Episoden mit Erfolgs-Label} lässt
    unbewertete weg.
  \item[2. Hardware] Gerät, Mischgenauigkeit, Datenlader.
  \item[3. Einstellungen] Schritte, Batchgröße, Zeitbudget; \knopf{Vorgaben wiederherstellen}.
  \item[4. Lauf] \knopf{Training starten}, \enquote{fortsetzen} für einen abgebrochenen Lauf,
    Fortschritt und Validierung.
\end{description}

\section{Reiter Betrieb}

\begin{description}
  \item[1. Modell] Checkpoint wählen. Die Oberfläche prüft, ob Kameras, Takt und Override zur
    Aufnahme passen.
  \item[2. Kameras] wie bei der Aufnahme; \knopf{Kameras prüfen} macht ein Probebild in der
    Fotostellung.
  \item[3. Fahrt] Anzahl Fahrten, Takte je Fahrt, \knopf{Inferenz starten}, \knopf{STOPP}.
    \enquote{Live-Ansicht} öffnet beim Start ein Fenster mit Bildern, \acs{tcp}-Spur und
    Vorhersagen.
  \item[4. Policy-Test] \enquote{Lage bekannt (Policy-Test gegen PICK)} mit
    \enquote{Abstand über Objekt} (Standard \SI{25}{\milli\metre}): Die Policy fährt an und wird
    über dem Bauteil angehalten (\cref{sec:policy-test}). \enquote{Experte greift und legt um}
    lässt danach den Planer greifen und an eine neue Lage legen; \enquote{DAgger-Ordner}
    speichert diese Fahrten.
\end{description}
